%-*-tex-*-
\ifundefined{writestatus} \input status \relax \fi %
\chcode{cont}

\def\cqu{\cquote{The life is so short, the craft so long to learn.}{Aphorisms,
Hippocrates (460-400~BC)}
}


\chapterhead{cont}{SPECIAL CHARACTERS\cr and\cr ACCENTS}
This chapter discusses {\it special characters}, both those that are single
keys (usually) on the terminal but are used for special purposes
and those, non mathematics,  that end up in the final text as
single characters. An example of the first is the |{| used for grouping (see
Chapter~\ref{bas} information on {\it grouping}) and an example of the latter is the copyright sign
\copyright\ that is obtained by typing |\copyright|. For special mathematics
symbols see Chapter~\ref{matsym}. 

Section~\ref{controluse} discusses the special characters that do not
normally print and are used by \tex\ and \intex\ for control purposes.
Section~\ref{controlprint} tells how to print them in the text.

A second set of special use characters, usually single keys, 
are really single character commands to \tex\
and \intex. Their behaviour varies depending on the circumstances. In most
cases, they can be used with little or no thought. Occasionally errors in use
might cause some confusion. For instance, in \intex\ the ``\vrt'' is
interpreted as a vertical bar or rule in normal text and mathematics, but as a
column separator with a vertical rule while building tables using \intex's
|\begintable...\endtable| commands. The complete set of these commands for
\intex\ is given in Section~\ref{activechar}.


The rest of this chapter discusses the second class of special {\it text}
characters, and the means of inserting them in the text.
Section~\ref{specialtextsymbols} shows how to obtain such interesting
characters as the British Pound ``{\it\$}'' paragraph symbol ``\P''. 
Quotes, both English ``\dots'' and French  
{\language = 1 \ldq\dots\rdq}  are discussed in
Section~\ref{quotes}.
Accents
on ordinary text letters are in Section~\ref{ordaccent} and on mathematics in
Section~\ref{mathaccent}. 



\shead{controluse}{Control Characters - Meaning}
\tex\ and \intex\ have reserved the following control characters for the uses
indicated. None of these characters will be actually printed in the text.
\bshortcomlist
|\|&escape character \dots starts 
                    all control and command sequences  \cr
|{}|&grouping brackets \cr
\% &indicates the end of a line in a \TeX\  file. Anything after the 
          \%\ will be ignored by \TeX\  while processing the file. This is used
       for comments. \cr
|\@|&superscript sign\cr
|_|&subscript sign\cr
|&|&alignment character used to align characters from line to line\cr
|$|&\it mathematics mode indicator\cr
|$<expression>$|&{\it mathematics} mode start and end \cr
|$$<expression>$$|&{\it display mathematics} mode \dots  start and end \cr
|#|&parameter indicator for definitions and in alignment preambles\cr
|~|&the {\it tie} character \dots\ really an unbreakable space between words.
For instance, the space between the initials  in D.E. Knuth allows the last
name to be split from the initials at the end of the line. The sequence 
|D.E.~Knuth| will leave the space but not allow the break.\cr
\eshortcomlist


\shead{controlprint}{Control Characters - Printing}
The control characters 
\begintt
      {     }     #     $     %     &      _     \     ~     ^
\endtt 
\ls
are obtained by the control sequences
\begintt
     $\{$    $\}$   \#    \$    \%    \&     \_     
        $\backslash$   \^{ }   \~{ }
\endtt
\ls 
respectively,  
except as accents. The last two are actually 
obtained by ``accenting'' a space. 
In mathematics the ``~\^{ }~'' and the ``~\~{ }~'' must be obtained using
the mathematic accents (see Section~\ref{mathaccent}) 
as |\hat{ }| or |\tilde{ }|. 
The net effect of this is that all characters are available in any mode. 

\shead{activechar}{\intex\ Special Characters}
\intex\  special characters are all active ---  meaning that they
are single character control sequences. The complete list is
given below. They are used for  specifying column separators
in tables.
\bshortcomlist
\vrt&this is a \vrt\ in normal text, a \vrt\ {\it Rel} in mathematics (see
Chapter~\ref{matsym}). While making a table inside \intex\
|\begintable ... \endtable|, the \vrt\ is a vertical rule or line  column
that is used to separate data columns and {\bf must 
not be used in mathematics inside tables -- use |\vrt|} 
    (see Chapter~\ref{align}).\cr
|"|&the double quote should not be used in ordinary text. \intex\ will give
you a gentle reminder to either |``| or |''|. It will then put the " in the
text. The |"| is used while making tables for indicating a rule column where
the rule is omitted.\cr
\eshortcomlist

\shead{specialtextsymbols}{Special Text Symbols}
\tex\ supplies a number of special symbols that can be used in ordinary text,
and some admonitions about those that should not. There are three kinds of
dashes in normal typeset material and \tex\ supplies all three. They are 
- -- --- and are produced by |-| |--| and |---| respectively. They are known
officially as the {\it dash, endash,} and {\it emdash}. They are used as
hyphens, word separators, and idea separators, respectively. 
These are recognized by \tex\ as {\it ligatures}. Although the English
ligatures such as ``ff'' and ``ffi'' are recognized by \tex, most in other
languages are not.\footnote{\dagger}{However, it is easy to modify \tex\ so
that it does recognize those ligatures. Meanwhile, it is necessary to live
with the mechanism given.}
These  require
special symbols. In addition there are special symbols such as 
\copyright\ and \dots\ that are also useful. 
Here is a list supplied by {\it plain}
\bshortcomlist
\@|\aa, \AA|& the \aa\ and \AA \cr
\@|\ae, \AE|& the ligature \ae\ or \AE \cr
\@|\copyright|&the copyright symbol \copyright \cr
\@|\dag|& the dagger \dag. The  math mode |$\dagger$| looks similar but 
                    incorporates the spacing of a binary operator. \cr
\@|\ddag|& the double dagger \ddag. The math mode
                    |$\ddagger$|  looks similar but incorporates the 
                     spacing of a binary operator. \cr
 |\dots|&dots in text \dots\ like so \cr
\@|{\it\$}|&the nonobvious British Pound {\it\$} \cr %indexing problem with $
\@|{\it\&}|&the script ampersand {\it\&} \cr 
\@|\i|&the dotless \i\ for use in \^\i \cr
\@|\j|&the dotless \j\ for use in \^\j \cr
\@|\l, \L|&the \l\ and \L \cr
\@|\o, \O|&the \o\ and \O \cr 
\@|\oe, \OE|&the ligature \oe\ or \OE \cr
\@|\P|&the paragraph symbol \P \cr
\@|\S|&the section symbol \S \cr
\@|\ss|&the \ss \cr
\eshortcomlist


\shead{quotes}{Opening and Closing Quotes}
The ``double quote'', |"| on the keyboard {\bf should not be used},
except in its accent form |\"| to be described later in
Section~{\ref{ordaccent}} on accents.
For the |\englishversion| of \intex, opening {\bf double} quotes may be obtained either
with the |``| in plain or using |\ldq|. The closing {\bf double} 
quotes may be obtained
similarly using the |''| or the |\rdq|. When |\versionfrancaise| is in force,
|\ldq| and |\rdq| when used like so  |\ldq guillemets\rdq|\footnote{\ddagger}{This is the best we can
do with the present fonts} will give {\language=1 \ldq guillemets\rdq}.
Note that the |\rdq| is an ordinary command that will delete spaces following
it. 


\shead{ordaccent}{Ordinary Accents}
The ordinary accents supplied in {\it plain} are for the occasional insertion
of a foreign language in  English text. The mechanisms are complete, and the
accents are quite well placed, but the method of obtaining them from
keyboards with special characters is not built in. In \tex, the accent always
precedes the letter accented. Thus |\'ecole| gives \'ecole, and |fran\c cais|
gives fran\c cais. Note the annoying space after the |\c| in the word. This
form of input reduces the readability of the text. Agreement is needed to
obtain a consistent multi-language format. 

The following is a list of the ordinary accents in \tex\ specified by {\it
plain}.
\bshortcomlist
|\`<char>|&accent grave as in \`a\cr
|\'<char>|&accent acute as in \'e \cr
|\^<char>|&circumflex or ``hat'' as in \^o \cr
|\"<char>|&umlaut as in \"o \cr
|\~<char>|&tilde as in \~n \cr
|\=<char>|&long vowel as in \=o \cr
|\.<char>|&dot accent as in \.a \cr
|\b|\]|<char>|&accent bar-under as in \b d, note space \] \cr
|\c|\]|<char>|&accent cedilla as in \c c, note space \] \cr
|\d|\]|<char>|&accent dot-under as in \d o, note space \] \cr
|\u|\]|<char>|&accent breve as in \u o, note space \] \cr
|\v|\]|<char>|&``check'' as in \v o \cr
|\H|\]|<char>|&long Hungarian umlaut as in \H o\cr
|\t|\]|<char><char>|&tie-after accent as in \t ao\cr 
\eshortcomlist


\shead{mathaccent}{Accents on Mathematical Symbols} 
Placing accents on mathematical symbols requires a certain amount of care. 
{\it Plain} \tex\ gives an extra set of commands that are to be used for characters in
mathematics mode. All of the following should be used when in mathematics
mode. Note the required space \], in almost all cases.
\begingroup \descrwd=2in
\bshortcomlist
|<char>'|& the prime as in $x'$ \cr
|<char>''|& the double prime as in $x''$ \cr 
|\acute|\]|<char>|&as in $\acute x$ \cr
|\bar|\]|<char>|&as in $\bar x$ \cr
|\breve|\]|<char>|&as in $\breve x$ \cr
|\check|\]|<char>|&as in $\check x$ \cr
|\dot|\]|<char>|&as in $\dot x$ \cr
|\ddot|\]|<char>|&as in $\ddot x$ \cr
|\grave|\]|<char>|&as in $\grave x$\cr
|\hat|\]|<char>|&as in $\hat x$ \cr
|\tilde|\]|<char>|&as in $\tilde x$ \cr
|\vec|\]|<char>|&as in $\vec x$ \cr
|\widehat|\]|<char>[<char>][<char>]|& as in $\widehat{x},\widehat{xy},
 \widehat{xyz}$\cr
|\widetilde|\]|<char>[<char>][<char>]|& as in $\widetilde{x},\widetilde{xy},
 \widetilde{xyz}$\cr
\eshortcomlist
\endgroup
When using these in a technical paper or report, it is probably better to
make a short definition for the accented character. This will reduce typing
and facilitate the changing the style {\t consistently} throughout the paper.
It is a very simple form of {generic typesetting}. 

\ejectpage
\done

